\documentclass[10pt,a4paper]{amsart}
\usepackage[margin=19mm]{geometry}
\usepackage{amsmath,amssymb,mathtools}
\usepackage[scr=rsfs,cal=euler]{mathalfa}
\usepackage{xfrac}
\usepackage[unicode,hidelinks,bookmarks=false]{hyperref}
\newcommand{\ie}{i.e.\ }
\newcommand{\cf}{cf.\ }
\newcommand{\resp}{respectively\ }
\newcommand{\Subsection}{\subsection}
\providecommand{\nobreakdash}{\nobreak\hbox{-}}
\setlength{\emergencystretch}{2em}
\multlinegap0pt
\usepackage{mathtools}
\usepackage[shortlabels]{enumitem}
\usepackage{xcolor}
\usepackage{csquotes}
\newtheorem{lemma}{Lemma}
\newcommand{\vp}{\nu_p}
\title{\bfseries Logarithm of the Universal\\ Two-Valued Formal Group}
\author{Victor Buchstaber, Mikhail Kornev}
\date{Source: arXiv:2609.15878v1 (CC BY 4.0)}
\begin{document}
\maketitle

\noindent\emph{Source and licence.} \bfseries Logarithm of the Universal\\ Two-Valued Formal Group. Version arXiv:2609.15878v1 (CC BY 4.0), distributed by arXiv under the Creative Commons Attribution 4.0 International licence, http://creativecommons.org/licenses/by/4.0/, as recorded in the arXiv licence metadata for this exact version and shown on the abstract page https://arxiv.org/abs/2609.15878v1. Full licence notice: \texttt{licenses/arxiv-2609.15878-CC-BY-4.0.txt}.\par\smallskip

\noindent\textit{Editorial scope.} Counted unit: the article's lemma lemma (source0.tex lines 2089--2101). The article's complete original proof of it is delivered with it (source0.tex lines 2103--2288, one \texttt{proof} environment of 186 lines). No mathematics is written, repaired or re-derived: the statement and the proof are the article's own lines, in the article's own notation, and no retained line is substituted or re-flowed.  No further source-local context is needed: every cross-reference of every retained line resolves inside the two delivered spans.  Nothing was omitted from any retained span, so the package carries no omission record.  No retained line contains a citation, so the wrapper carries no bibliography and no bibliography entry.  No retained line contains a figure environment, an image-inclusion command or drawing code, so no image is delivered and the package declares no asset dependency.  Editorial formatting only, in addition: an AMS \texttt{amsart} preamble that restates the article's own declarations and macros used by the retained lines, namely the environments \texttt{lemma} on separate counters, together with the standard packages those lines need.  The retained environments print in this document as Lemma 1 in order of appearance, whereas the article prints the same items with the numbering of its own full document.  The complete native source of the article as distributed in the arXiv e-print for this exact version is bundled unchanged as \texttt{sources/210458/source0.tex}.
\begin{lemma}[(Arithmetic lemma)]\label[lemma]{lem:arithmetic}
For \(m\ge2\), put
\[
\alpha_{m,k}=d_{m-1}r_{m,k},
\qquad 1\le k\le m.
\]
Then $\alpha_{m,k}\in\mathbb Z$ for all \(k\), and
\[
\gcd(\alpha_{m,1},\ldots,\alpha_{m,m})=1.
\]
Thus the numbers \(r_{m,k}\) have the common denominator \(d_{m-1}\)
with relatively prime numerators.
\end{lemma}

\begin{proof}
Put \(D=d_{m-1}\), so that
\[
\alpha_{m,k}=Dr_{m,k}.
\]
We shall prove the stronger identity
\begin{equation}\label{eq:max-valuation}
\max_{1\leq k\leq m}\bigl(-\vp(r_{m,k})\bigr)=\vp(D)
\end{equation}
for every prime \(p\).

The identity
\begin{equation}\label{eq:r-simplified}
r_{m,k}
=
\frac{1}{m}
\frac{\binom{m}{k}}
{k\binom{2k-1}{k-1}}
\end{equation}
gives
\begin{equation}\label{eq:valuation-r}
-\vp(r_{m,k})
=
\vp(m)
+
\vp\!\left(k\binom{2k-1}{k-1}\right)
-
\vp\binom{m}{k}.
\end{equation}

First let \(p\) be odd and put
\[
e=\left\lfloor\log_p(2m-1)\right\rfloor.
\]
Because $p$ is odd and $2m$ cannot be a power of $p$,
\begin{equation}\label{eq:odd-D}
\vp(D)=\vp(m)+e.
\end{equation}
By Legendre's formula,
\begin{align}
\vp\!\left(k\binom{2k-1}{k-1}\right)
&=
\vp\!\left(\frac{(2k-1)!}{(k-1)!^2}\right)\notag\\
&=
\sum_{j\geq1}
\left(
\left\lfloor\frac{2k-1}{p^j}\right\rfloor
-
2\left\lfloor\frac{k-1}{p^j}\right\rfloor
\right).
\label{eq:legendre-sum}
\end{align}
If \(k-1=Ap^j+s\), where \(0\leq s<p^j\), then the \(j\)-th
summand equals
\[
\left\lfloor\frac{2s+1}{p^j}\right\rfloor.
\]
It is either \(0\) or \(1\), and it vanishes when \(p^j>2k-1\).
Consequently,
\[
\vp\!\left(k\binom{2k-1}{k-1}\right)\leq e,
\]
and \eqref{eq:valuation-r} gives
\begin{equation}\label{ineq}
-\vp(r_{m,k})\leq\vp(m)+e=\vp(D).
\end{equation}

We now show that equality in \eqref{ineq} is attained. Put \(t=p^e\).
If \(t\leq m\), take \(k=t\). Lucas' theorem gives
\[
\vp\binom{m}{k}=0,
\]
while \(k-1=p^e-1\), so every summand in
\eqref{eq:legendre-sum}, for \(1\leq j\leq e\), equals \(1\).
Hence equality holds in \eqref{ineq}.

Suppose that \(m<t\). Then
\[
\frac{t+1}{2}\leq m<t.
\]
Put \(h=(p-1)/2\) and \(q=(t+1)/2\). In base \(p\), the digits of
\(q\) are
\[
q_0=h+1,
\qquad
q_i=h\quad (1\leq i<e).
\]
Write \(m=\sum_{i=0}^{e-1}m_ip^i\). If \(m=q\), set \(\ell=0\).
Otherwise, let \(\ell\) be the largest index for which \(m_\ell\neq q_\ell\).
In the latter case, $m>q$ and hence $m_\ell>q_\ell$. In the former case,
$m_0=q_0=h+1$. Define \(k\) by the base-\(p\)
digits
\[
k_i=
\begin{cases}
0,&i<\ell,\\
h+1,&i=\ell,\\
h,&i>\ell.
\end{cases}
\]
Every digit of \(k\) is at most the corresponding digit of \(m\), so
Lucas' theorem gives
\[
\vp\binom{m}{k}=0.
\]
The digits of \(k-1\) are equal to \(p-1\) below position \(\ell\) and
to \(h\) from position \(\ell\) onwards. Therefore, for every
\(1\leq j\leq e\),
\[
(k-1)\bmod p^j\geq\frac{p^j-1}{2},
\]
and hence every summand in \eqref{eq:legendre-sum} equals \(1\).
Thus
\[
\vp\!\left(k\binom{2k-1}{k-1}\right)=e,
\]
so equality in \eqref{ineq} is again attained.

It remains to consider \(p=2\). Legendre's formula gives
\begin{equation}\label{eq:binary-weight}
\nu_2\!\left(k\binom{2k-1}{k-1}\right)=s_2(k-1),
\end{equation}
where \(s_2(r)\) denotes the number of ones in the binary expansion of
\(r\).

If \(m=2^e\), then \(s_2(k-1)\leq e\) for every \(k\leq m\), and
\eqref{eq:valuation-r} gives
\[
-\nu_2(r_{m,k})\leq\nu_2(m)+e=\nu_2(D).
\]
For \(k=m\),
\[
s_2(k-1)=e,
\qquad
\binom{m}{k}=1,
\]
so equality is attained.

Suppose that \(m\) is not a power of \(2\), and put
\[
e=\left\lceil\log_2m\right\rceil.
\]
Then \(s_2(k-1)\leq e-1\) for every \(k\leq m\), and therefore
\[
-\nu_2(r_{m,k})
\leq
\nu_2(m)+e-1
=
\nu_2(D).
\]
Taking \(k=2^{e-1}\), Lucas' theorem gives
\[
\nu_2\binom{m}{k}=0,
\qquad
s_2(k-1)=e-1,
\]
so equality is attained in this case as well. This proves
\eqref{eq:max-valuation} for every prime \(p\).

It follows from \eqref{eq:max-valuation} that
\[
\vp(\alpha_{m,k})
=
\vp(D)+\vp(r_{m,k})
\geq0
\]
for every prime \(p\) and every \(k\). Hence
\[
\alpha_{m,k}\in\mathbb Z.
\]

Finally, for each prime \(p\), choose an index \(k\) for which the maximum
in \eqref{eq:max-valuation} is attained. Then
\[
\vp(\alpha_{m,k})
=
\vp(D)+\vp(r_{m,k})
=
0.
\]
Thus no prime divides all the integers \(\alpha_{m,1},\ldots,\alpha_{m,m}\).
Consequently,
\[
\gcd(\alpha_{m,1},\ldots,\alpha_{m,m})=1.
\]
\end{proof}

\end{document}
